\ContentFor{ADD}{

\section{\unclass  Introduction}  
%%%%% If you are making changes to this tex file, please speak to technical leadership. 

\guidance{	\unclass TIMELINE: Should be finalized prior to SRR, and unless it's a brand new type of system, should be reasonably understood prior to kickoff.
	\unclass This section provides a high-level overview of the document’s purpose and scope. It defines the intended audiences–-such as users, developers, and stakeholders–-and explains how the document supports their understanding
and use of the model. It establishes the foundation for all subsequent sections by outlining the intended use for
the model and the structure of the documentation.

} % ends guidance

\classlabelshort The \modelname{} \gls{add} discusses the codebase for the \modelname{} model developed by \gls{bah} in the \gls{odessa} framework in support of the \gls{ngic}.  The information in this document outlines the computational decisions, frameworks, and development pathways that guided the development of this model. This document prepares users to understand the codebase at a much deeper computer science level, whereas the UAM is concerned more with user guidance, the underlying theoretical background and mathematics, and only broad-view discussion of the codebase. 

} % ends \ContentFor{ADD}


\ContentFor{UAM}{
	
	\section{\unclass  Introduction}  
	%%%%% If you are making changes to this tex file, please speak to technical leadership. 
	
	\guidance{	\unclass TIMELINE: Should be finalized prior to SRR, and unless it's a brand new type of system, should be reasonably understood prior to kickoff.
		\unclass This section provides a high-level overview of the document’s purpose and scope. It defines the intended audiences–-such as users, developers, and stakeholders–-and explains how the document supports their understanding
		and use of the model. It establishes the foundation for all subsequent sections by outlining the intended use for
		the model and the structure of the documentation.
		
	\unclass EXAMPLE: This document is a companion to the primary documentation on the Overarching Dynamic Electronic-Warfare System Standard Architecture (ODESSA) architecture which specifically outlines the relevant data on the \modelname{} \textit{type} model. The information in this document outlines the equations, mathematical algorithms, theoretical background, and gives context to the decisions and methodologies which guided the development of the EXAMPLE MODEL NAME model. This document prepares users and analysts to work with the EXAMPLE MODEL NAME by detailing its parameters and providing a setup guide.
		
	} % ends guidance
	
\classlabelshort The \modelname{} \gls{uam} discusses the purpose, scope, and function of the \modelname{} \textit{type} model developed by \gls{bah} in the \gls{odessa} framework in support of the \gls{ngic}.

} % ends \ContentFor{UAM}



%%%%%%%%%%%%%%%%%%%%%%%%%%%%%%%%%%%%%%%%%%%%%

\ContentFor{TDD}{

\section{\unclass  Introduction}  

\unclass This document summarizes the test plan for the \modelname{} \gls{tmap} model developed by \gls{bah} in the \textcolor{red5}{[architecture name here]} (ex: \gls{odessa}) for \gls{ngic}. The purpose of this \gls{tdd} is to define a strategy for verification of the model implementation. Verification is the process of determining the model and associated data accurately represent the conceptual description and specification as described in the model requirements. Validation is the determination that the model is an accurate representation of the real world from the perspective of the intended use and will be performed by \gls{ngic}.

\classlabelshort The \modelname{} \gls{tdd} documents the testing suite for the \modelname{} \textit{type} model developed by \gls{bah} in the \gls{odessa} framework in support of the \gls{ngic}.

} % ends \ContentFor{TDD}

%%%%%%%%%%%%%%%%%%%%%%%%%%%%%%%%%%%%%%%%%%%%

\ContentFor{ICD}{

\section{\unclass  Introduction}  
%%%%% If you are making changes to this tex file, please speak to technical leadership.

\guidance{\unclass TIMELINE: This information should be filled out by PDR. \\

\unclass For each described DLL specify which version of the architecture is associated with it.

\unclass If your model is not being wrapped for integration, this document can simply state that the model is
only designed to interface with the ODESSA Scenario GUI and currently does not support external integration. \\

} 
	
\classlabelshort This document describes the use of the compiled \gls{dll} of the X.Y.Z version of the \gls{odessa} framework.  \gls{bah} creates models for use with \gls{odessa} in support of the \gls{ngic} to facilitate integration of the model and \gls{odessa} into external simulation environments.

} % ends \ContentFor{ICD}

%%%%%%%%%%%%%%%%%%%%%%%%%%%%%%%%%%%%%%%%%%%%%